% =====================================================================
\section{실험환경}
\label{sec:env}

\subsection{문헌 기반 합성환경}
\label{sec:env-calib}
제안 아키텍처는 특정 터미널의 운영기록이 아니라 문헌 기반 합성환경에서 평가한다.
환경의 범위는 \ref{sec:related}절의 선행연구를 참고했다. 여기에는 20개 블록
~\cite{ref16}, 블록별 복수 크레인~\cite{ref8}, 여러 주의 연속 운영~\cite{ref9},
트럭 도착의 뚜렷한 시간대별 변동~\cite{ref13}, 기항량이 다른 세 선박 등급
~\cite{ref16}, 크레인 취급시간 범위~\cite{ref9}가 포함된다. 이 범위 안에서 본 연구는
21개 블록, 블록당 크레인 2기, 상태가 이어지는 30일 연속 운영, 세 선박 등급과 180초의
기준 크레인 서비스시간을 사용한다. 도착곡선은 \ref{sec:env-demand}절에서 설명한다. 이
값들은 본 연구의 설계 선택이다. 인용문헌은 값이 기존 모형의 범위 안에 있음을
뒷받침하지만, 특정 항만의 운영자료에서 직접 추정했다는 뜻은 아니다.

\subsection{터미널과 시간 설정}
\label{sec:env-terminal}
야드는 21개 블록과 42기의 야드 크레인으로 구성한다. 각 블록에는 크레인 2기를 두고,
용량은 1,440슬롯, 초기 점유율은 65\%로 설정한다. 게이트와 안벽은 1차원 야드의 서로
반대쪽에 있으며 결정 간격은 60초다. 운영은 30일 동안 연속으로 진행하고 가운데 28일을
비교한다. 크레인은 실제 서비스 작업을 먼저 처리하고, 그 안에서는 예상 처리시간이
짧은 작업을 우선한다. 하루의 경계는 비용을 구분하는 시점일 뿐 터미널을 초기화하는
시점이 아니다. 야드 재고, 대기열, 크레인 상태와 미완료 작업은 다음 날로 이어진다.
첫날과 마지막날은 이 연속 경계를 연결하는 데 사용한다.

\subsection{외부 트럭 작업 수요와 도착 과정}
\label{sec:env-demand}
외부 트럭 작업 수요는 선행연구가 보고한 범위를 참고해 합성했다. 선행연구는 비슷한
수요 규모와 하루 중 큰 도착 변동을 뒷받침한다~\cite{ref13}. 다만 아래의 구체적인
도착곡선은 본 연구가 구성한 것이다. 일일 수요는
$\{3{,}500,\,5{,}000,\,7{,}500,\,12{,}500,\,15{,}000\}$ 대 위의 이산분포에서 확률
$\{0.30,\,0.30,\,0.25,\,0.10,\,0.05\}$로 추첨한다(그림~\ref{fig:demand}a). 앞의 세
수준은 원활, 보통과 높은 수요를 나타내고, 뒤의 두 수준은 혼잡과 초혼잡 조건을
나타낸다. 각 작업의 통지시각은 도착밀도 $f(t)$에서 생성한다. 이 밀도는 전체 물량의
38\%를 담는 균등 기저와 나머지 62\%를 담는 세 개의 가우시안 성분으로 구성한다
(그림~\ref{fig:demand}b).

\begin{equation}
f(t)=0.38\,U(0,24)+0.62\sum_{m=1}^{3}w_m
\mathcal N(t;\mu_m,\sigma_m^2),
\end{equation}

여기서 $(\mu_m,\sigma_m,w_m)$은 각각 $(10,1.5,.317)$, $(15,2.5,.633)$와
$(21,1,.05)$이고 시간 단위는 시(hour)다. 균등 기저는 야간 도착률을 0보다 크게
유지한다. 따라서 날짜 경계를 넘어 이어지는 야드 상태가 인위적으로 비어 있는 밤의
결과가 되지 않는다. 봉우리의 위치, 폭과 가중치는 수요수준별 혼잡 민감도를 비교하기
위한 설계값이며 특정 터미널에서 추정한 값이 아니다. \textbf{학습과 평가는 같은
곡선 형태를 사용한다.} 두 단계는 수요 수준과 난수만 다르다. 따라서 형태가 다른
도착곡선에 대한 일반화는 평가하지 않았다.

\begin{center}
\begin{minipage}{\textwidth}
\centering
\includegraphics[width=\textwidth]{fig-demand.pdf}
\captionof{figure}{외부 트럭 작업 수요. (a) 다섯 개의 일일 수요 수준과 각각이 뽑힐 확률이며
날마다 독립으로 추첨한다. (b) 하루 12{,}500대 조건의 시간대별 도착 과정으로, 균등
기저와 가우시안 성분 셋을 이들이 합쳐 이루는 곡선 아래에 따로 그렸다. 선행연구가
뒷받침하는 것은 수요 규모와 도착이 하루 중 균일하지 않다는 사실이고, 혼합의 형태는
본 연구의 구성이다.}
\label{fig:demand}
\end{minipage}
\end{center}

\subsection{선박과 안벽 작업}
\label{sec:env-vessel}
선박은 소형, 중형과 대형의 세 등급으로 나눈다. 표~\ref{tab:vessels}는 등급별 규모,
배정 STS, 유휴비용과 일일 기항 구성을 요약한다.

\begin{center}
\captionof{table}{선박 등급과 일일 기항 구성. 기항 비율은 하루 평균 3.3척을 기준으로
계산했다. 대형 0.3척은 매일 30\% 확률로 1척을 추가하는 설정의 기대값이다.}
\label{tab:vessels}
\small
\setlength{\tabcolsep}{4pt}
\begin{tabular}{@{}lrrrrr@{}}
\toprule
등급 & GT & TEU & STS & 유휴비용 & 일일 기항 \\
 & & & (기) & (천원/h) & (평균 척수 비율) \\
\midrule
소형 & 50,000 & 3,000 & 2 & 149.5 & 2.0 (60.6\%) \\
중형 & 100,000 & 7,500 & 4 & 299.0 & 1.0 (30.3\%) \\
대형 & 150,000 & 14,000 & 6 & 448.5 & 0.3 (9.1\%) \\
\bottomrule
\end{tabular}
\end{center}

용량이 $C$ TEU인 선박의 기항 작업량은 $u\sim U[0.15,0.30]$에 대해 $C\cdot u$로
추출한다. STS 1기의 처리율은 27.5\,move/h이며, 이는 문헌의 24, 30과
36\,move/h 범위 안에 있다~\cite{ref16}. 배정된 STS에서 발생한 막힘은 합산한 뒤
STS 수로 나누어 선박별 유휴시간으로 환산한다.

야드 크레인 1기의 기준 서비스시간은 180초/move이므로 이론상 총 이동능력은
$\mu_{YC}=21\times2\times3600/180=840$ move/h다. move는 컨테이너 이동 횟수이며
트럭 건수와 다르다. 트럭 작업 하나가 여러 이동과 재취급을 요구할 수 있으므로 두 단위를
구분해 보고한다.

\subsection{비용함수}
\label{sec:env-cost}
날짜 구간 $[d,d+1)$의 총비용은 네 항의 합으로 정의한다.

\begin{equation}
\Phi_d=C_{wait,d}+C_{move,d}+C_{rehandle,d}+C_{vessel,d}.
\end{equation}

트럭 체류비는 시간당 40,000원이며, 1시간을 넘는 구간에는 같은 단가를 한 번 더
적용한다. 화폐 기준점은 안전운임에서 가져왔고 시간 환산 방식은 본 연구의 설계값이다
~\cite{ref24}. 추가 크레인 이동은 작업시간당 60,000원, 재취급은 회당 2,000원으로
설정한다. 두 값도 연구 설계값이다. 선박 유휴비는 접안 초과요율 29.9원/10GT·시간을
GT·시간당 2.99원으로 환산했다~\cite{ref25}. 트럭 $i$의 체류시간을 $h_i$, 추가 크레인
작업시간을 $m$, 재취급 횟수를 $R$, 선박 $v$의 정책 관련 유휴시간을 $T_v$라 하면
비용은 다음과 같다.

\begin{equation}
\begin{aligned}
C_{wait}&=\sum_i 40{,}000\{h_i+\max(0,h_i-1)\},\qquad
C_{move}=60{,}000\,m,\\
C_{rehandle}&=2{,}000\,R,\qquad
C_{vessel}=\sum_v 2.99\,GT_v T_v.
\end{aligned}
\end{equation}

접이안과 검역처럼 재배치 정책이 바꿀 수 없는 구조적 선박 시간은 별도 진단값으로
분리한다. 네 비용 항은 같은 날짜 경계에서 측정한 누적값의 차이로 계산한다.

학습에 사용한 30일 실행에서 환경의 두 가지 성질을 확인했다. 첫째, 대기비용이
총비용을 지배한다. $C_{wait}$는 96.4\%를 차지하며, 재취급은 2.2\%, 선박 유휴는
1.2\%, 추가 크레인 이동은 0.1\%다. 따라서 $\Phi$에 대한 판단은 대부분 규정에서
가져온 안전운임 계수에 의해 결정된다~\cite{ref24}. 문헌값이 아닌 두 설계 계수의 영향은
상대적으로 작다. 둘째, 수요가 증가하면 대기비중과 서비스 수준이 급격히 나빠진다.
대기비중은 하루 3{,}500대에서 83.2\%, 7{,}500대에서 91.7\%, 15{,}000대에서
99.1\%로 증가한다. 한 달 전체의 평균 체류시간은 0.89시간이고 90분위 체류시간은
2.01시간이다. 트럭 작업 206{,}904건 중 5.0\%는 통지된 날 안에 끝나지 않았다. 특히
90분위 체류시간은 가장 한산한 날 0.79시간에서 가장 붐비는 날 7.49시간까지 9.5배
증가했다. 즉 혼잡한 날은 비용 규모뿐 아니라 비용 구성도 다르다. 재배치 정책은 바로 이
차이를 활용해야 한다.
